\documentclass[11pt]{article}
\makeatletter\def\input@path{{../estilo/}}\makeatother
\usepackage{didistribuidores}
\renewcommand{\documentocorto}{Bitácora · Etapa 2}

\begin{document}
\portada{Bitácora de pruebas}{Etapa 2: notificaciones como servicio aparte}%
  {Outbox, RabbitMQ con cola persistente y de mensajes fallidos, y un consumidor idempotente}%
  {Rama \cod{02-separacion}}{Noviembre de 2026}

\tableofcontents
\newpage

\section{Qué se construyó en esta etapa}

En la etapa 2, \cod{notificaciones} sale del núcleo y pasa a ser un servicio con \textbf{su propio
proceso, su propia base de datos} (\cod{notificaciones\_db}) y su propia API (en esta etapa, en el
puerto 8001). El núcleo ya no lo llama: cada cambio de estado de un pedido deja un evento en
\cod{pedidos.outbox}, un publicador lo manda a una cola persistente de RabbitMQ y el servicio de
notificaciones lo consume y guarda el aviso.

\begin{center}
\begin{tikzpicture}[node distance=0.9cm and 1.2cm, every node/.style={font=\sffamily\small},
  caja/.style={draw=azulsd, fill=grisfondo, rounded corners=3pt, minimum height=0.9cm, align=center, text width=2.9cm},
  fl/.style={-Stealth, thick, azulsd}]
  \node[caja] (nuc) {Núcleo\\\scriptsize cambio de estado + evento en \texttt{outbox} (1 transacción)};
  \node[caja, right=of nuc] (pub) {Publicador\\\scriptsize hilo del núcleo, con confirmación};
  \node[caja, right=of pub] (rmq) {RabbitMQ\\\scriptsize cola durable \texttt{notificaciones.pedidos}};
  \node[caja, below=of rmq] (con) {Consumidor\\\scriptsize valida, guarda y luego \emph{ack}};
  \node[caja, left=of con] (db) {\texttt{notificaciones\_db}\\\scriptsize \texttt{evento\_id} único};
  \node[caja, below=of nuc, draw=rojo] (dlq) {Cola de fallidos\\\scriptsize \texttt{...pedidos.dlq}};
  \draw[fl] (nuc) -- (pub); \draw[fl] (pub) -- (rmq); \draw[fl] (rmq) -- (con); \draw[fl] (con) -- (db);
  \draw[fl, rojo] (con.south) |- ++(0,-0.6) -| node[pos=0.25,below,font=\sffamily\scriptsize]{nack sin reencolar} (dlq.south);
\end{tikzpicture}
\end{center}

El camino de un aviso:
\begin{enumerate}
  \item El pedido cambia de estado y, \textbf{en la misma transacción}, se guarda el evento en el
        \cod{outbox}. En ese momento el pedido ya está confirmado: nadie espera a ningún aviso.
  \item Un publicador (un hilo del núcleo) lee los eventos pendientes y los publica en RabbitMQ. Solo
        los marca como publicados cuando RabbitMQ confirma que los recibió y los guardó
        (\emph{publisher confirms}).
  \item Los mensajes esperan en una cola durable y son persistentes: sobreviven aunque el consumidor
        esté apagado o RabbitMQ se reinicie.
  \item El consumidor valida cada mensaje, guarda el aviso y \textbf{solo después} hace \emph{ack}.
  \item Si el mismo evento llega dos veces, se ignora (\cod{evento\_id} es \cod{UNIQUE}); si un mensaje
        no se puede interpretar, va a la cola de mensajes fallidos y no bloquea a los demás.
\end{enumerate}

\begin{nota}[Al menos una vez + idempotente]
No se puede garantizar que un mensaje llegue \emph{exactamente} una vez entre dos procesos que pueden
caerse en cualquier momento. Sí se garantiza que llega \textbf{al menos una vez} (outbox, confirmación,
cola durable y \emph{ack} tardío) y que repetirlo \textbf{no tiene efecto} (consumidor idempotente).
Juntas, las dos garantías equivalen a ``exactamente una vez'' para el cliente.
\end{nota}

Cada evento lleva la \textbf{versión del pedido}, que sube con cada cambio de estado. Como los mensajes
pueden llegar desordenados, es ese número, y no la fecha, el que dice cuál es el estado más reciente:
los relojes de dos máquinas nunca coinciden exactamente.

\subsection{Criterios del PRD que se comprueban (historia 5.2)}
\begin{tabularx}{\linewidth}{@{}l X l@{}}
\toprule
\textsf{\textbf{Criterio}} & \textsf{\textbf{Qué dice}} & \textsf{\textbf{Prueba}} \\
\midrule
CA1 & Con \cod{notificaciones} apagado, un pedido se crea, se paga y cambia de estado normalmente & 2 \\
CA2 & Al volver a encenderlo, llegan los avisos pendientes, cada uno una sola vez & 2 y 3 \\
CA3 & Un mensaje inválido va a una cola de mensajes fallidos y no bloquea a los demás & 4 \\
\bottomrule
\end{tabularx}

\section{Levantar el sistema}

\begin{comando}
docker compose down -v
docker compose up --build -d
docker compose ps
curl.exe -s http://localhost:8000/salud
curl.exe -s http://localhost:8001/salud
docker compose logs app | Select-String "Publicador"
docker compose logs notificaciones | Select-String "Consumidor"
\end{comando}

El \cod{down -v} es necesario porque \cod{db/04-notificaciones.sql}, que crea
\cod{notificaciones\_db}, solo se ejecuta cuando PostgreSQL arranca con el volumen vacío. Estas
pruebas se hicieron en PowerShell de Windows (\cod{curl.exe}, \cod{Invoke-RestMethod},
\cod{Select-String}); en Linux los comandos equivalentes usan \cod{curl} y \cod{grep}.

\observado Cuatro servicios: la base y RabbitMQ en \cod{healthy}, la aplicación y notificaciones en
\cod{Up}. Las dos rutas de salud respondieron \cod{\{"estado":"ok"\}} y
\cod{\{"estado":"ok","consumidor\_conectado":true\}}. En los registros aparecen
\emph{Publicador del outbox conectado a RabbitMQ} y
\emph{Consumidor de notificaciones esperando eventos en notificaciones.pedidos}.

\figura{img/fig-01.png}{Los cuatro servicios, las rutas de salud y los registros del publicador y del consumidor.}

\subsection{Consola de RabbitMQ}
En \url{http://localhost:15672} (usuario \cod{sd}, contraseña \cod{sd}), pestaña de colas, existen la
cola principal \cod{notificaciones.pedidos} y la de mensajes fallidos \cod{notificaciones.pedidos.dlq}.
Las dos tienen la marca \textbf{D} (\emph{durable}: sobreviven a un reinicio del broker). La principal
muestra además \textbf{DLX} y \textbf{DLK}: tiene configurado a dónde mandar los mensajes rechazados.
Al arrancar, las dos están en cero.
\figura{img/fig-02.png}{Consola de RabbitMQ: la cola principal y la de fallidos, ambas durables y vacías.}

\section{Preparación de las pruebas}

Se creó un usuario de la zona 1, se tomaron un restaurante y un repartidor de la misma zona, y se
definieron cuatro funciones auxiliares en PowerShell:

\begin{tabularx}{\linewidth}{@{}l X@{}}
\toprule
\textsf{\textbf{Función}} & \textsf{\textbf{Qué hace}} \\
\midrule
\cod{Ciclo-Completo} & Crea un pedido y lo lleva hasta \cod{ENTREGADO} (pagar, aceptar, entregar con el código). Devuelve su id. \\
\cod{Avisos} & Consulta los avisos de un pedido en el servicio de notificaciones (versión, tipo y mensaje). \\
\cod{Cola} & Muestra los mensajes y consumidores de cada cola (\cod{rabbitmqctl list\_queues}). \\
\cod{Psql} & Ejecuta una consulta en la base del núcleo. \\
\bottomrule
\end{tabularx}

\section{Prueba 1: flujo normal}

\queria Que con todo encendido, cada cambio de estado de un pedido termina como un aviso en el
servicio de notificaciones.

\begin{comando}
$id = Ciclo-Completo; Start-Sleep 3
Avisos $id
Psql "SELECT tipo, publicado_en IS NOT NULL AS publicado FROM pedidos.outbox
      WHERE payload->>'pedidoId'='$id' ORDER BY creado_en"
\end{comando}

\esperado Tres avisos (versiones 2, 3 y 4) y tres eventos con \cod{publicado = t}.

\observado Tres avisos: \emph{Recibimos tu pago de \$30.00. Estamos buscando un repartidor.},
\emph{Un repartidor aceptó tu pedido.} y \emph{Tu pedido fue entregado. ¡Buen provecho!}. En el
\cod{outbox}, los tres eventos del pedido con \cod{publicado = t}: salieron de la tabla, pasaron por
RabbitMQ y llegaron a la base de notificaciones.
\figura{img/fig-03.png}{Flujo normal: tres avisos y tres eventos publicados.}

\section{Prueba 2: notificaciones apagado (criterios 1 y 2)}

\queria Que los pedidos no dependen del servicio de avisos y que los avisos no se pierden mientras
está apagado.

\begin{comando}
docker compose stop notificaciones
$id = Ciclo-Completo
(Invoke-RestMethod "$base/api/pedidos/$id").estado
Start-Sleep 3; Cola
\end{comando}

\observado Con notificaciones detenido, el pedido llegó de todas formas a \cod{ENTREGADO} (criterio 1).
La cola principal tiene \textbf{3 mensajes y 0 consumidores}; la de fallidos está vacía.
\figura{img/fig-04.png}{Con notificaciones apagado el pedido llega a \cod{ENTREGADO} y la cola acumula 3 mensajes.}

\begin{comando}
docker compose start notificaciones
Start-Sleep 6
Avisos $id
Cola
\end{comando}

\observado Al encenderlo llegaron los tres avisos, en orden y una sola vez cada uno (criterio 2), y la
cola quedó en 0 mensajes con 1 consumidor.
\figura{img/fig-05.png}{Al volver notificaciones llegan los tres avisos y la cola queda vacía.}

\section{Prueba 3: evento duplicado}

\queria Que el consumidor es idempotente: aunque RabbitMQ entregue el mismo mensaje más de una vez, el
aviso se guarda una sola vez.

\comofunciona La herramienta \cod{herramientas/publicar\_prueba.py} publica eventos de prueba
directamente en RabbitMQ; con \cod{--evento-id} y \cod{--veces} se publica el mismo evento varias veces.

\begin{comando}
$antes = (docker compose logs notificaciones | Select-String "Evento repetido").Count
$PD = [guid]::NewGuid(); $EV = [guid]::NewGuid()
docker compose exec -T notificaciones python herramientas/publicar_prueba.py `
    --pedido-id $PD --evento-id $EV --veces 3
Start-Sleep 2; Avisos $PD
$despues = (docker compose logs notificaciones | Select-String "Evento repetido").Count
"duplicados ignorados: $($despues - $antes)"
\end{comando}

\observado Un solo aviso guardado y \textbf{2 duplicados ignorados}. \cod{evento\_id} tiene \cod{UNIQUE}
en \cod{notificaciones\_db} y el \cod{INSERT} usa \cod{ON CONFLICT DO NOTHING}.
\figura{img/fig-06.png}{Un evento publicado tres veces: un aviso guardado y dos duplicados ignorados.}

\section{Prueba 4: mensaje inválido (criterio 3)}

\queria Que un mensaje que no se puede procesar va a la cola de fallidos y no bloquea a los demás.

\comofunciona El consumidor distingue dos tipos de error. Los \textbf{permanentes} (JSON mal formado, tipo
desconocido, falta un campo) se rechazan con \cod{nack} sin reencolar y RabbitMQ los desvía a la cola
de fallidos: reintentarlos nunca funcionaría. Los \textbf{temporales} (la base no responde) cierran la
conexión sin \emph{ack} y el mensaje se reentrega al reconectar.

\begin{comando}
docker compose exec -T notificaciones python herramientas/publicar_prueba.py --invalido
Start-Sleep 2; Cola
$PD = [guid]::NewGuid()
docker compose exec -T notificaciones python herramientas/publicar_prueba.py `
    --pedido-id $PD --tipo pedido.entregado --version 4
Start-Sleep 2; Avisos $PD
\end{comando}

\observado Un mensaje en \cod{notificaciones.pedidos.dlq} y la cola principal en cero. El evento válido
publicado enseguida se guardó con normalidad: el mensaje roto no bloqueó la cola.
\figura{img/fig-07.png}{Mensaje inválido en la cola de fallidos y evento válido posterior guardado.}

\section{Prueba 5: eventos desordenados}

\queria Que si los eventos de un pedido llegan fuera de orden, el sistema sabe cuál es el más reciente.

\begin{comando}
$PD = [guid]::NewGuid(); $UD = [guid]::NewGuid()
docker compose exec -T notificaciones python herramientas/publicar_prueba.py `
    --pedido-id $PD --usuario-id $UD --tipo pedido.entregado --version 4
docker compose exec -T notificaciones python herramientas/publicar_prueba.py `
    --pedido-id $PD --usuario-id $UD --tipo pedido.asignado --version 3
Start-Sleep 2; Avisos $PD
Invoke-RestMethod "http://localhost:8001/api/notificaciones/pedido/$PD/ultimo" | Format-List tipo, version_pedido
\end{comando}

\observado La versión 4 se publicó antes que la 3. Aun así, los avisos salen ordenados por versión (3 y
luego 4) y \cod{/ultimo} responde \cod{pedido.entregado} con versión 4: el estado más reciente, aunque
haya llegado primero.
\figura{img/fig-08.png}{Eventos desordenados: avisos ordenados por versión y \cod{/ultimo} con la versión 4.}

\section{Prueba 6: RabbitMQ caído}

\queria Que los pedidos siguen funcionando con RabbitMQ apagado y que los eventos no se pierden: esperan
en el \cod{outbox} hasta que el broker vuelva.

\begin{comando}
docker compose stop rabbitmq
$id = Ciclo-Completo
(Invoke-RestMethod "$base/api/pedidos/$id").estado
Psql "SELECT count(*) AS pendientes FROM pedidos.outbox WHERE publicado_en IS NULL"
\end{comando}

\observado El pedido avanzó hasta \cod{ENTREGADO} y el \cod{outbox} mostró \textbf{3 eventos pendientes}:
se guardaron en la misma transacción que el cambio de estado.
\figura{img/fig-09.png}{RabbitMQ apagado: el pedido llega a \cod{ENTREGADO} y sus 3 eventos quedan pendientes.}

\begin{comando}
docker compose start rabbitmq
Start-Sleep 40
Psql "SELECT count(*) AS pendientes FROM pedidos.outbox WHERE publicado_en IS NULL"
Avisos $id
\end{comando}

\observado Cuarenta segundos después, los pendientes bajaron a 0 y llegaron los tres avisos. El
publicador y el consumidor se reconectaron solos.
\figura{img/fig-10.png}{RabbitMQ de vuelta: 0 pendientes y los tres avisos del pedido.}

\section{Prueba 7: la cola sobrevive a un reinicio de RabbitMQ}

\queria Que los mensajes que esperan en la cola no se pierden si RabbitMQ se reinicia.

\begin{comando}
docker compose stop notificaciones
$id = Ciclo-Completo; Start-Sleep 3; Cola
docker compose restart rabbitmq; Start-Sleep 30; Cola
docker compose start notificaciones; Start-Sleep 8; Avisos $id
\end{comando}

\observado La cola tenía 3 mensajes antes y después del reinicio, y al encender notificaciones llegaron
los tres avisos. El mensaje que se ve en la cola de fallidos es el inválido de la prueba 4, que también
sobrevivió.

\begin{nota}[Por qué RabbitMQ tiene \cod{hostname} fijo]
RabbitMQ guarda sus datos según el nombre del nodo. Sin un \cod{hostname} fijo en el Compose, cada
contenedor nuevo tendría otro nombre y las colas persistentes parecerían desaparecer al recrearlo.
\end{nota}
\figura{img/fig-11.png}{Los mensajes sobreviven al reinicio de RabbitMQ y los avisos llegan completos.}

\section{Resumen de la etapa 2}

\begin{longtable}{@{}p{4.3cm} p{7.2cm} l l@{}}
\toprule
\textsf{\textbf{Prueba}} & \textsf{\textbf{Resultado observado}} & \textsf{\textbf{Fig.}} & \textsf{\textbf{Estado}} \\
\midrule
Levantar los cuatro servicios & \cod{db} y \cod{rabbitmq} sanos; consumidor conectado; colas durables creadas & 1--2 & \ok \\
1. Flujo normal & Tres avisos (versiones 2, 3 y 4) y tres eventos publicados & 3 & \ok \\
2. Notificaciones apagado & El pedido llega a \cod{ENTREGADO}; 3 mensajes esperan; al volver llegan una vez cada uno & 4--5 & \ok \\
3. Evento duplicado & Un aviso; 2 duplicados ignorados & 6 & \ok \\
4. Mensaje inválido & Un mensaje en la cola de fallidos; el siguiente evento se guarda & 7 & \ok \\
5. Eventos desordenados & Avisos ordenados por versión; \cod{/ultimo} da la versión 4 & 8 & \ok \\
6. RabbitMQ caído & El pedido avanza; 3 pendientes; al volver, 0 pendientes y llegan los avisos & 9--10 & \ok \\
7. Reinicio de RabbitMQ & 3 mensajes antes y después; avisos completos & 11 & \ok \\
\bottomrule
\end{longtable}

\subsection*{Comprobado otra vez con el código final}
La regresión del documento de pruebas (sección de la etapa 2) se volvió a correr el 18 de noviembre
de 2026 sobre el código con las cuatro etapas, ya por el gateway: los tres avisos de un pedido
completo llegan con versiones 2, 3 y 4, los tres eventos del \cod{outbox} quedan publicados y la cola
principal queda en 0 mensajes con 1 consumidor.

\section{Quién hizo qué}

\begin{tabularx}{\linewidth}{@{}p{4.2cm} X@{}}
\toprule
\textsf{\textbf{Integrante}} & \textsf{\textbf{Aportes en la etapa 2}} \\
\midrule
Yukioayax Canek Gabriel Hernández (@YukiGab) & Servicio de notificaciones (base propia, consumidor, API y cola de mensajes fallidos), publicador del \cod{outbox} hacia RabbitMQ, topología compartida, RabbitMQ en el Compose, herramienta \cod{publicar\_prueba.py} y las siete pruebas de la etapa. \\
Erick Isaac Echeverria Goicochea (@ErickI07) & Revisión del pull request de la etapa 2. \\
Los dos & Decisiones de alcance y de diseño del PRD; contrato de eventos (\cod{docs/eventos.md}); revisión cruzada; evidencias y bitácora. \\
\bottomrule
\end{tabularx}

\end{document}
